% ===========================================================================
%  English translation of CN.docx
%  Source: CN.docx (135,103 bytes, sha256 8cb27428...). The source is unchanged.
%  Translation is faithful: numbers, claims and hedges are carried over as written.
%  Discrepancies found in the source are listed in the translator's notes at the
%  end rather than silently corrected.
% ===========================================================================
\documentclass[12pt,a4paper]{article}

\usepackage[T1]{fontenc}
\usepackage[utf8]{inputenc}
\usepackage{lmodern}
\usepackage[margin=2.5cm]{geometry}
\usepackage{amsmath,amssymb}
\usepackage{booktabs}
\usepackage{longtable}
\usepackage{array}
\usepackage{enumitem}
\usepackage{microtype}
% --- typographic scale ------------------------------------------------------
% 12 pt with 1.5x line spacing. Chosen by measurement, not by taste: at 11 pt
% single spacing the literature review occupies 4 pages, and the target was 7-8.
% The same size and spacing apply to every chapter, since both are set in the
% preamble. Content is unchanged -- only these two parameters differ.
\usepackage{setspace}
\setstretch{1.5}
\usepackage[hidelinks]{hyperref}

\newcommand{\qt}{q_t}
\newcommand{\wt}{\omega_t}

\title{Independent Reproduction and Review of LMCompress across Four Modalities}
\author{}
\date{}

\begin{document}
\maketitle

% ---------------------------------------------------------------------------
\begin{abstract}
\noindent
The goal of this report is to independently implement and review the compression
pipeline of LMCompress in four modalities: text, audio, image and video. In the course
of the implementation I rewrote the arithmetic coder, the measurement engine and the
data loaders, and designed a unified calling interface for the three models Llama, bGPT
and ImageGPT. Throughout this reproduction I used the authors' released training weights
and test sets so that the reproduction could be checked against them; the authors'
released code was used only for reference and cross-checking.

For the text base model, Llama, I sampled from the MeDAL test set; the sampled scale is
smaller than the full evaluation in the original paper. The base model reaches a rate of
\textbf{9.1353\%} under an 8{,}192-character window, a difference of \textbf{0.08
percentage points} from the 9.06\% reported in the paper. The LoRA adapter released by the
authors has an effect in the \emph{opposite} direction to that claimed in the paper: it
raises the rate under all five settings examined, with a penalty of
\textbf{0.58--0.68 percentage points}.

For the audio model bGPT I ran a same-payload test on the LJSpeech test set. On the same
82{,}220-byte audio clip, the independent encoder emits 16{,}297.875\,B against the
paper's reference value of 16{,}399.95\,B, a relative difference in compression rate of
about \textbf{0.1\%}.

For the image and video model ImageGPT, the image modality could not be brought into
agreement with the paper, and it could not be shown that the original paper's procedure is
lossless; the video modality shares the same problem of a non-lossless compression
procedure. For images, the median rate over five CLIC images is \textbf{21.4781\%},
roughly seven percentage points from the paper's 14.62\%. For video, the per-channel test
over nine distinct frames gives a rate of \textbf{26.1866\%}, also markedly higher than the
paper. Neither images nor video undergo a lossless procedure: iGPT encodes through a
non-injective preprocessing step that resizes to $32\times32$ and maps colours onto 512
palette entries, so the original pixels are not what is coded, and without side information
no lossless reconstruction is possible. An exhaustive check shows that the 256 possible
single-channel values map onto only \textbf{48} distinct token IDs, so the compression
procedure supplied by the source code is not lossless.

This work shows that part of LMCompress can be reproduced independently: the text base rate
and the audio result are close to the published values. The claimed benefit of LoRA for the
text model could not be reproduced, and the image and video results are not sufficient to
support the paper's conclusion that the procedure is lossless. The conclusions of the
original paper need to be read modality by modality.
\end{abstract}

% ---------------------------------------------------------------------------
\section{Introduction}

\subsection{Background}

The requirement of lossless compression is unambiguous: the decompressed data must be
identical to the data before compression. Lossless compression applies wherever information
loss cannot be accepted---medical records and images, legal documents, scientific data,
archival storage. What a lossless compressor must do is exploit as much of the predictable
part of the data as possible in order to reduce the number of bits needed to represent it,
while retaining all the information required to recover the original data exactly.

Shannon's information theory shows that the more predictable the data, the shorter the
average code length needed to describe it. In the development of classical lossless
compressors a compression system came to be split into two cooperating stages: a
\emph{predictor} that estimates the probability of the next symbol from the content already
processed, and an \emph{encoder} that uses those probabilities to assign a code length to
the symbol that actually occurs. The quality of the predictor determines how closely the
probabilities match the true distribution of the data, while the encoder is responsible for
turning probabilities into a bitstream that can be transmitted and decoded.

Large pre-trained models give the predictor stage a new means of realisation. The
statistical regularities a large model learns from massive data can be used to predict the
probability of the next symbol in a stream of text, audio or visual information. But accurate
prediction does not by itself mean that the system achieves lossless compression of the
source data: how the data are turned into tokens, whether the decoder can obtain the same
predictive distribution, and whether the preprocessing retains the original information all
determine whether exact reconstruction is ultimately possible.

\subsection{Object of study and research questions}

This project takes as its object of study the reproduction of LMCompress, proposed by Li
\emph{et al.} That paper combines autoregressive models suited to different modalities with
arithmetic coding. Its core framework is to convert the input into a token sequence, predict
the next token from the known context, and then encode the actual sequence using the
predicted probabilities. The paper applies this framework to four modalities---text, audio,
image and video---using a different model and data representation for each, and compares
against the classical compression baselines predominant in that modality.

This reproduction of LMCompress is concerned not only with compression rate but also with
whether the procedure the paper proposes constitutes a complete lossless compression system.
The reproduction therefore addresses the following questions:

\begin{enumerate}[label=\arabic*.,leftmargin=*]
  \item Can the main pipeline be implemented independently, without calling the authors'
        compression routines, and can the paper's results be reproduced under the same or
        similar settings?
  \item What representation does the model actually encode, and does it correspond to the
        source data used as the denominator in the paper? Can the original input be
        recovered after decoding?
  \item If reproduced values differ from those in the paper or in the source code, can the
        difference be explained by specific factors such as window length, sample, scoring
        position or data preprocessing?
\end{enumerate}

\subsection{Scope and work performed}

This reproduction implements the text, audio, image and video pipelines independently and
compares local results against both the formally published paper and the authors' reference
implementation. For text, the compression rate of the base model and of the model with a
LoRA adapter attached are examined separately, with character windows and token windows
kept explicitly distinct, and measurements completed under several conditions. For audio,
the compression rate was reproduced while also checking the correspondence between the byte
payload, the patches, and the output positions of the model. For image and video the
per-channel input method proposed in the paper is followed, and the question of whether
iGPT's palette quantisation makes the compression procedure lossless is examined further.
Each modality is accompanied by a corresponding classical compression baseline, describing
the behaviour of traditional methods on the same input.

Three lines of work were emphasised. First, the interface between the large-model predictor
and the arithmetic coder was completed independently, so that predictions become a code
length that can actually be checked. Second, the measurement basis was inspected along the
whole path from data representation to coded output, paying particular attention to whether
the symbols the model actually processes and the denominator used in the report refer to the
same object. Third, where results differed, the respective experimental conditions were
re-established so that the causes of the difference could be sorted out and explained, with
results compared under identical samples and windows.

On this basis the report distinguishes three levels: the model's probabilistic prediction
over the token sequence, the bitstream the encoder emits, and the ability to recover the
original data from that bitstream. Only when the decoder can reproduce the prediction
process, decode the tokens accurately, and the preprocessing representation retains all the
information of the source data does a low rate constitute evidence of lossless compression
of the source data.

% ---------------------------------------------------------------------------
\section{Literature review}

\subsection{Scope of the review}

This review follows one central question through the development of lossless compression
research: how can a system use the information it already has in order to represent the next
symbol of the source data in fewer bits, while retaining the ability to restore the source
information exactly? Early classical compressors achieved compression by counting repeated
substrings within a window; later compressors performed prediction using hand-designed
context models. In the new period in which neural networks and large language models serve
as predictors, greater model scale and predictive capability, context-window settings, and
application across modalities have become new questions within the same line of development.

At the same time, because lossless compression is the premise, the object obtained after
decompression must be exactly the object supplied at compression. If the decoder can only
restore a transformed sequence, and that transformation is many-to-one, then the original
data have not been recovered. A low model loss does not by itself amount to a compressed
file either: the probabilities the model supplies must be converted into a code that a
decoder with the same conditions can parse correctly. Understanding LMCompress therefore
requires connecting source-coding theory, prediction, practical entropy coding,
representation design and experimental measurement.

This chapter follows that line: information-theoretic bounds and coding methods first, then
dictionary compression and context-based compression, then learned probabilistic models and
large pre-trained models, then how images and audio are converted into symbol sequences, and
finally LMCompress together with the conventions needed to evaluate a reproduction.

\subsection{Discrete entropy, entropy rate and coding}

Shannon defined the discrete entropy of a discrete random variable through its probability
mass function. This expression fixes the ``lower bound'' of source coding: under lossless
coding the average code length cannot fall below the entropy of the source.

For sources with memory, the limit of lossless compression is determined by the entropy
rate. Real data are usually not drawn independently from a fixed distribution; the
probability of a symbol typically depends on context---surrounding bytes, preceding words,
neighbouring pixels, or earlier samples in an audio stream. The chain rule allows the
probability of a sequence to be written as a product of conditional probabilities. Put into
practice for compression, this means the predictor estimates the probability of the next
symbol from the symbols already obtained.

\paragraph{Encoders.}
Given the chain rule and Shannon's theory, the more accurate the predicted probabilities,
the shorter the code length. The encoder's task is precisely to map the probability
distribution supplied by the predictor onto a uniquely decodable code satisfying the
Kraft--McMillan inequality, keeping each symbol's code length as close as possible to its
information content: symbols with high probability should receive shorter codewords, and
symbols with low probability may use longer ones. Following this idea, Huffman and Witten
proposed Huffman coding and arithmetic coding in 1952 and 1987 respectively. The former
gives an optimal prefix-code construction for a distribution over a finite alphabet,
reaching a global optimum through locally optimal choices, and was used in early classical
lossless compressors. Huffman coding builds an optimal prefix code for a known finite
alphabet and suits symbol-by-symbol coding; it is simple and efficient, but each codeword
must be an integral number of bits. When a symbol's probability approaches 1, an ordinary
binary prefix code still needs at least one bit, and when many such symbols are coded
separately the Huffman code length can be appreciably above the ideal information content.
Arithmetic coding removes the one-bit-per-symbol floor: it does not allocate a separate
codeword to each symbol but successively narrows an interval according to the probabilities
of consecutive symbols. The width of the interval representing the whole sequence depends on
the probability of that sequence, and a single binary number inside the interval identifies
it. Engineering realisations such as range coding and binary arithmetic coding are widely
used in later classical lossless compressors and in large-model-based lossless compressors.

\subsection{Theoretical and practical realisation of classical compressors}

\paragraph{The Lempel--Ziv line.}
Lempel--Ziv methods match repeated substrings that have previously occurred in the context;
the DEFLATE specification is the typical representative. DEFLATE was proposed in 1981 and
combines LZ77, proposed in 1977, with Huffman coding. LZ77 uses a sliding window to count
past repeated fragments, maintaining a table that assigns a number to every (literal,
length) pair; dynamically counting the frequencies of those numbers yields the probability
distribution of the file. DEFLATE then hands that dynamic distribution to a Huffman tree.

Dictionary methods are particularly effective when the data contain repeated patterns: they
need not estimate a complete probability distribution over all possible files. The matching
rule also limits their scope, however. They readily identify exactly repeated content but
are poor at capturing semantic similarity or long-range relationships that do not manifest
as repeated strings. Window size, block boundaries and dictionary-reset policy all affect
how much history a compressor can exploit.

\paragraph{Universal modelling and coding.}
LZ77 compresses while it predicts, so combining LZ77 with Huffman coding does not separate
predictor from encoder. In the work considered here the compression task is for the first
time decomposed into two independent modules, prediction and coding. From then on the design
of the predictor became a research subject in its own right.

Separating prediction from coding is very helpful for research. Replacing the predictor does
not require redesigning the arithmetic coder; the same coder can be paired with a
hand-built context model, a recurrent network or a Transformer. The division also gives a
clearer way to compare models: first report the code length implied by the model's
probabilities, then the actual coded output, and explain where the two differ.

\paragraph{Variable-order context models.}
Once predictor and encoder were decoupled, attention shifted to designing the predictor.
Variable-order context models in the Markov tradition became an important line, with LOEMA,
PPM and PAQ as typical representatives.

LOEMA (Local Order Estimating Markovian Analysis) was the first lossless compressor to use a
variable-order Markov model as its predictor, predicting the next symbol's probability with
several orders simultaneously and using the confidence of each order to decide the most
suitable order locally. The confidence design relies on hand engineering, however, and
predicting with several orders at once brings low speed and high computational cost.

PPM estimates the next symbol's probability using the longest-order Markov model, and when a
symbol has never been seen it encodes an ESC (escape) and falls back to the next shorter
order. This solves the cost of predicting with several models simultaneously, but there is no
uniform answer to how much probability the escape symbol should receive; different escape
designs gave different PPM variants such as PPM-A, PPM-B, PPM-C and PPM-D. PPM's fixed
maximum order also prevents it from capturing longer-range dependencies.

The earliest version of PAQ resembles LOEMA in idea, retaining the predictive power of
several Markov orders simultaneously, but instead of choosing the most suitable order it
\emph{mixes} the predictions. PAQ's compression performance depends heavily on hand-designed
implementation details.

\subsection{Neural networks and modern large-model compressors}

The classical methods above trace the development of the compression system and point to how
the predictor should evolve: any replacement must still assign valid probabilities to every
possible next symbol using the context, and produce a bitstream through reversible coding.
With the rapid development of neural networks, predictors built on neural networks entered
researchers' field of view.

\paragraph{DeepZip and Dzip.}
DeepZip combines a recurrent neural network with arithmetic coding and was evaluated on text,
genomes and synthetic sequences. The work also observes that the average code length per
symbol agrees closely with the cross-entropy, indicating that the cross-entropy loss used to
train a model is essentially the same objective as the minimum code length sought in
compression.

Subsequent neural compressors explored how to train or adapt the predictor for the file being
compressed. DZip further combines neural modelling with adaptive and semi-adaptive training,
attempting to reduce dependence on a large external training corpus. Such methods bring
practical trade-offs: more adaptation to the input may improve probability estimates, but
model state, training time, memory footprint and the decoder's obligation to reproduce the
adaptation process all become part of the system design. A method that yields a short ideal
code length may still, if compressing a medium-sized file takes days, be better suited to
research than to everyday storage.

\paragraph{Lossless data compression with a Transformer.}
The paper by Izacard, Joulin and Grave was the first to extend the neural approach to
Transformers, still using arithmetic coding as the encoder and evaluating on enwik8. Its
results show that a purely neural predictor is competitive for lossless compression---that a
Transformer can serve as an effective lossless-compression predictor. The route from lossless
compression to large models was thereby fully opened.

\paragraph{LLMZip.}
Having established the effectiveness of Transformer predictors for lossless compression,
research on large-model-based lossless compression concentrated on whether pre-trained models
can serve as transferable compression priors, and LLMZip is an important node on that path.
In the text domain it combines LLaMA-7B with arithmetic coding and achieves high compression
rates across several text datasets. Its main contribution is to show that a pre-trained model
can predict probabilities effectively.

\paragraph{LMCompress.}
LMCompress, proposed by Li \emph{et al.}, extends the large-model-predictor-plus-arithmetic-
coding approach from text to images, video and audio. The report applies three models to
these four modalities, with a workflow that stays consistent throughout: convert the raw data
into a token sequence suitable for the model, have the autoregressive model estimate the
conditional probability of the next token, and perform entropy coding with arithmetic coding.
According to the reported results, LMCompress achieves better compression rates than the
classical baselines on the datasets chosen for all four modalities, and breaks earlier records
for lossless compression with large models.

The contribution of this work lies in its breadth: it tests whether one general idea extends
from natural language to other data types. Because several modalities are covered, the raw
image and audio data must first be converted into tokens the model can process, and how that
conversion is done necessarily affects experimental results and conclusions. Likewise, the
text tokenisation scheme, the audio patch layout, the image processor and the video frame
sampling rule jointly define what the model actually evaluates.

\subsection{Datasets}

LMCompress's multi-modal evaluation covers images, video, audio and text.

For images it uses ILSVRC and CLIC2019 Professional, concentrating on the processing of
natural images. For video it uses Xiph.org Video Test Media (Derf's Collection), which
provides a variety of uncompressed YUV4MPEG test sequences spanning different resolutions,
frame rates and scene content. These visual datasets supply uncompressed sequences, avoiding
the influence of prior compression.

For audio LMCompress uses LibriSpeech and LJSpeech. The former consists of multi-speaker
English read speech from LibriVox public-domain audiobooks; the latter is read speech from a
single female speaker. Combining multi-speaker and single-speaker data lets LMCompress test
compression behaviour under different speech regularities.

For text LMCompress uses MeDAL and Pile of Law. MeDAL contains large numbers of medical
abbreviations and specialised expressions, while Pile of Law collects legal and administrative
text. By testing on medical and legal text---domains where lossless compression and
transmission are required---LMCompress demonstrates good compression capability in practical
application domains.

\subsection{Summary}

The literature above shows an increasingly close connection between artificial intelligence
and information compression. Early work grew out of Shannon's theory and encoders, and in
exploring the compression of real text it gradually produced the decoupling of predictor and
encoder with research focused on the predictor; the relationship between prediction and
textual context gave rise to several predictors designed in the Markov tradition. After the
rise of neural networks, research using neural networks and Transformers as predictors
extended into the new field of large-model-based lossless compression.

LMCompress extends this direction to several modalities, compressing text, images, audio and
video with large models. Its results suggest that the semantic understanding learned by large
models can fundamentally improve compression efficiency and surpass traditional hand-designed
methods.

Reproducing LMCompress therefore has both technical significance and considerable academic
value.

% ---------------------------------------------------------------------------
\section{Method}

\subsection{Clean-room implementation}

The implementation of LMCompress's compression logic was carried out independently. The
authors' code was used in exactly four ways: (1) as a reference object for reading the
logic; (2) as the source of pre-trained weights; (3) as a supplementary source of data
reference; and (4) as an unmodified executable, in order to study its runtime behaviour.

\subsection{Definition of the measurement basis}

The important quantities used in this report are defined as follows.

\begin{enumerate}[label=\arabic*.,leftmargin=*]
  \item $N$: the number of symbols counted in the payload. For images and video this follows
        the number of processed samples as defined in the paper.
  \item $\qt$: the token probability distribution the model gives at position $t$.
  \item $\wt$: the token that actually occurs at position $t$.
  \item $C$: the sum, in bits, obtained by converting the probability of each actual token
        under the model distribution into a code length:
        \begin{equation}
          C \;=\; \sum_{t=1}^{N} -\log_2 \qt(\wt).
        \end{equation}
\end{enumerate}

The compression rate follows the original paper, computed as compressed bits divided by the
bits of the original payload, each payload symbol counting as 8 bits when uncompressed:
\begin{equation}
  \text{rate} \;=\; \frac{C}{8N}.
\end{equation}
The compression ratio is the reciprocal of the rate:
\begin{equation}
  \text{ratio} \;=\; \frac{1}{\text{rate}} \;=\; \frac{8N}{C}.
\end{equation}

It should be noted that $C$ is the model-implied code length computed from the model's
float32 probabilities, not the final coded length of the arithmetic coder. In
implementation the arithmetic coder quantises probabilities into integer CDFs, so the final
output differs slightly.

\subsection{Pipeline design}

Following the paper and the structure of the source code, data in all four modalities pass
through three stages:

\begin{enumerate}[label=\arabic*.,leftmargin=*]
  \item \textbf{Tokenisation.} Raw data are converted into token sequences the model can
        process. Text goes through Llama's BPE tokeniser; audio fills bytes into fixed-size
        patches; images and video are first split by colour channel and then quantised and
        tokenised by the ImageGPT processor.
  \item \textbf{Prediction.} Given the token sequence, the autoregressive model for the
        modality produces, at each position $t$, the distribution $\qt$ over the next token.
  \item \textbf{Coding.} An engineered arithmetic coder encodes the token sequence according
        to $\qt$.
\end{enumerate}

\begin{table}[htbp]\centering\small
\caption{Models and the working points reported in the paper.}
\begin{tabular}{@{}llll@{}}
\toprule
Modality & Model & Tokenisation & Reported setting \\
\midrule
Text & Llama-3-8B $\pm$ LoRA & BPE (Llama-3 vocabulary) & character windows 4{,}096/6{,}144/8{,}192; token windows 4{,}096/8{,}192 \\
Audio & bGPT (byte level) & byte $\rightarrow$ 16-byte patch & mono 8-bit PCM; prediction per patch \\
Image & iGPT-small & $32\times32$, 512-colour palette & single-channel patch replicated to greyscale RGB; 1{,}024 tokens per block \\
Video & iGPT-small (per frame) & as image, per frame & per frame, per channel; no inter-frame information \\
\bottomrule
\end{tabular}
\end{table}

In what follows, LMCompress$^{*}$ denotes the text model without fine-tuning, and
LMCompress the text model with the LoRA adapter attached.

\subsection{Integer range coder}

Our arithmetic coder uses fixed-point integer arithmetic to subdivide the unit interval. The
coder maps the whole probability range onto the integer range $[0, 2^{\text{precision}})$ and
then codes using the interval corresponding to the symbol. For an interval with bounds
$\mathrm{lo}$ and $\mathrm{hi}$, writing the current width as $\mathrm{width} = \mathrm{hi} -
\mathrm{lo}$, the new interval is
\begin{align}
  \mathrm{lo}' &= \mathrm{lo} + \left\lfloor \frac{\mathrm{width}\cdot \mathrm{cdf}[s]}{F} \right\rfloor, \\
  \mathrm{hi}' &= \mathrm{lo} + \left\lfloor \frac{\mathrm{width}\cdot \mathrm{cdf}[s+1]}{F} \right\rfloor,
\end{align}
where $F$ is the full integer range and $\mathrm{cdf}$ the integer cumulative distribution
over the alphabet.

The implementation was checked for basic correctness: every symbol is assigned a non-zero
integer interval, guaranteeing that a very small probability is not rounded to zero under
finite precision. The interval update is checked before output, and an exact round-trip test
was performed to confirm that the coding interval remains valid after every update and that
the symbol sequence recovered matches the input exactly.

\subsection{Text tokenisation and windows}

Text is tokenised with the BPE tokeniser of Llama-3-8B. The paired evaluation uses three
character windows (4{,}096, 6{,}144, 8{,}192) and two token windows (4{,}096, 8{,}192); the
two window units are reported separately. Each target token at position $t$ is paired with
the causal distribution at position $t-1$, and the independent code checks that targets and
token positions are aligned.

\subsection{Audio byte patches}

Following the original design, audio is processed as a byte sequence. bGPT uses 19 bytes per
patch and returns 17 slot predictions per patch, including the predicted first byte of the
next patch. The reproduction computes the probability the model assigns to each actual byte
using the corresponding patch and slot index. Where a patch cannot be filled, the reference
implementation pads the file to a multiple of 16 bytes so that every patch is full, whereas
our implementation fills the final incomplete patch with the special value 256. Although the
padding differs, it has no effect on the resulting compression rate.

\subsection{Image and video}

Images are predicted with iGPT-small. As described in the paper, each step processes a single
colour channel of the original image, first cutting that channel into $32\times32$ patches.
To make the single-channel data conform to the RGB three-channel format iGPT expects, the
code replicates each pixel value three times, i.e.\ a grey pixel whose R, G and B values are
equal. This yields 1{,}024 quantised tokens. Video reuses this per-channel path frame by
frame.

The official image path computes the code length of the token at position $t$ from the output
distribution at position $t$, and because that distribution already contains information about
the current token it does not satisfy the causal order required for step-by-step decoding. In
the independent implementation I changed this to predict the token at position $t$ from the
output distribution at position $t-1$.

\subsection{Classical compressor baselines}

I also reproduced the classical compressors used for comparison in the paper, including
OptimFROG, WavPack, FLAC and FFmpeg, tested on the datasets supplied with the paper. The
results show that these compressors handle the files correctly and that the computed rates are
within a reasonable range and similar to the paper's, indicating that the rate computation is
sound.

\subsection{Paper results used as the reference target}

The original paper's results, and where they appear, are given here as the reference
standard. ``CR'' in the paper denotes the compression rate (compressed/original, lower is
better); the reciprocal (ratio) is also tabulated.

It should be added that the preprint and the formally published \emph{Nature Machine
Intelligence} version differ in MeDAL sample size and in the fine-tuning results, so all
results here follow the formally published supplementary information (SI).

\begin{table}[htbp]\centering\small
\caption{Results taken from the paper as the reference standard.}
\begin{tabular}{@{}llrl@{}}
\toprule
Source & Item & Rate (\%) & Ratio ($\times$) \\
\midrule
SI Table 6 & MeDAL, LMCompress$^{*}$ (base) & 9.08 & 11.01 \\
SI Table 6 & MeDAL, LMCompress (LoRA) & 9.06 & 11.04 \\
SI Table 6 & Pile of Law, LMCompress$^{*}$ & 5.86 & 17.06 \\
SI Table 6 & Pile of Law, LMCompress (LoRA) & 4.90 & 20.41 \\
SI Table 6 & MeDAL, zpaq / brotli / 7z / zstd & 25.11 / 27.93 / 30.25 / 30.45 & --- \\
SI Table 3 & CLIC2019, LMCompress-igpt & 14.62 & 6.84 \\
SI Table 3 & CLIC2019, PNG / WebP / JPEG-XL & 47.08 / 35.56 / 31.16 & --- \\
SI Table 4 & Video high-resolution, LMCompress-igpt & 15.31 & 6.53 \\
SI Table 4 & Video high-resolution, H.265 / H.264 / FFV1 / zstd & 21.69 / 37.42 / 40.16 / 51.04 & --- \\
SI Table 9 & MeDAL at context 8{,}192 / 6{,}144 / 4{,}096 & 9.06 / 9.16 / 9.36 & --- \\
\bottomrule
\end{tabular}
\end{table}

\begin{table}[htbp]\centering\small
\caption{What was measured, and on what data.}
\begin{tabular}{@{}p{1.6cm}p{3.4cm}p{3.6cm}p{5.0cm}@{}}
\toprule
Modality & Data & Scale & Note \\
\midrule
Text & MeDAL sample (3 files) & 312{,}053\,B payload, 57{,}547 tokens & paired window evaluation; the paper uses a larger corpus \\
Audio & LJ001-0004.wav & 82{,}220\,B, 16\,kHz mono 8-bit PCM & single clip; the paper's result comes from its LJSpeech evaluation \\
Image & CLIC professional validation (5 images) & 40{,}648{,}704 counted source samples & the model processes per-channel patches; corpus scale differs from the paper \\
Video & Netflix Aerial 4K clip & 9 distinct frames; 26{,}542{,}080 samples per frame & per-channel model evaluation; classical baselines cover 10 frames of the same clip \\
\bottomrule
\end{tabular}
\end{table}

\subsection{How the reproduction was verified}

The principal values listed in this report can be re-checked against the raw measurements in
the appendix and the scripts listed there. Self-checks specific to the independent
implementation were also performed: whether the encoder decodes back to the source data,
whether the text prediction positions are aligned, whether the mapping from palette pixels to
tokens behaves as expected, and whether the independent results agree with the reference
values.

% ---------------------------------------------------------------------------
\section{Results and analysis}

\subsection{Overview}

\begin{table}[htbp]\centering\small
\caption{Summary of results by modality.}
\begin{tabular}{@{}p{1.5cm}p{2.5cm}p{4.2cm}p{2.4cm}p{3.0cm}@{}}
\toprule
Modality & Object & Current result & Reproduction verdict & Paper / note \\
\midrule
Text & base Llama-3-8B & A100 paired: 8.8038\,\%--9.4240\,\%; 9.1353\,\% at 8{,}192 characters; independent port 9.0410\,\% & close to the paper & SI Table 9: 9.06\,\% (+0.08\,pp); SI Table 6: 9.08\,\% ($-$0.04\,pp) \\
Text & authors' released LoRA & rate rises under all five settings; +0.58--0.68\,pp & direction differs from the paper & SI Table 6: 9.08\,\% $\rightarrow$ 9.06\,\% ($-$0.02\,pp) \\
Text baseline & brotli / zstd / 7z / zlib & fixed-window results in the baseline table of this chapter; brotli q11 gives 31.84\,\% at 8{,}192\,B & close to the paper's window grid & windows 8{,}192/16{,}384/32{,}768\,B; unbounded listed separately \\
Audio & bGPT & port 19.82\,\%; reference implementation 19.94\,\% & same payload, close & paper LJSpeech: 18.26\,\% \\
Image & ImageGPT causal path & CLIC five-image weighted 24.9892\,\%; median 21.4781\,\%; 20.7429\,\% excluding the highest & source image not recoverable & palette quantisation is many-to-one \\
Image baseline & WebP / JPEG-2000 / PNG & 39.36 / 50.22 / 76.85\,\% & classical codec results & 12 CLIC images \\
Video & ImageGPT, per frame, per channel & 9 distinct frames, 26.1866\,\% & rate of the quantised representation & paper 15.31\,\%; locally 10.88\,pp higher \\
Video baseline & H.265 / H.264 / FFV1 / zstd & 26.47 / 34.09 / 40.11 / 46.29\,\% & source-video coding baseline & 10 frames of the same clip \\
\bottomrule
\end{tabular}
\end{table}

\subsection{Text}

Three MeDAL files, 312{,}053\,B of payload in total. The paired evaluation covers the
4{,}096/6{,}144/8{,}192 character windows and the 4{,}096/8{,}192 token windows.

\begin{table}[htbp]\centering\small
\caption{Text rates by window unit and length.}
\begin{tabular}{@{}lrrrr@{}}
\toprule
Window unit and length & Base rate (\%) & +LoRA rate (\%) & Change (pp) & Effective context \\
\midrule
characters: 4{,}096 & 6.5322 & 6.9319 & +0.3997 & about 739 tokens/window \\
characters: 6{,}144 & 6.3927 & 6.7958 & +0.4031 & about 1{,}130 tokens/window \\
characters: 8{,}192 & 6.3321 & 6.7308 & +0.3987 & about 1{,}477 tokens/window \\
tokens: 4{,}096 & 6.1657 & 6.6107 & +0.4450 & about 3{,}837 tokens/window \\
tokens: 8{,}192 & 6.1023 & 6.5751 & +0.4728 & about 6{,}394 tokens/window \\
\bottomrule
\end{tabular}
\end{table}

The results for the five paired window settings are as follows.

\begin{table}[htbp]\centering\small
\caption{Text results across the five paired window settings.}
\begin{tabular}{@{}p{4.2cm}p{3.2cm}p{3.4cm}p{2.8cm}@{}}
\toprule
Statistic & Local result & Paper / reference result & Note \\
\midrule
Base model, five window settings & 8.8038\,\%--9.4240\,\% & 8{,}192 characters: 9.1353\,\%; independent port: 9.0410\,\% & window units listed separately \\
+LoRA, five window settings & 9.4859\,\%--10.0006\,\% & 8{,}192 characters: 9.5154\,\% (reference script) & paired measurement within the same corpus \\
Adapter change in rate & +0.5752 to +0.6821 percentage points & +0.5901\,pp (reference script) & positive under all five settings \\
Directional consistency & rate rises under all five settings & same direction as the reference script & local sample is three files \\
\bottomrule
\end{tabular}
\end{table}

On the three MeDAL files the compression rate of the base model LMCompress$^{*}$ varies with
the window setting over 8.8038\,\%--9.4240\,\%. The result under the 8{,}192-character
window, which corresponds to Supplementary Table 9 of the paper, is 9.1353\,\%, about
0.05 percentage points from the paper's 9.08\,\%.

\subsubsection{Evaluation of the LoRA adapter}

In the data source used here---SI Table 6 of the formal publication---fine-tuning on MeDAL
gives LMCompress$^{*}$ a rate of 9.08\,\% and LMCompress with the LoRA adapter 9.06\,\%, a
gain of 0.02\,pp or 0.2\,\% relative. On Pile of Law the reported result improves from
5.86\,\% for the base model to 4.90\,\%, a relative gain of 16\,\%.

The training parameters given in the paper are: 64\,MB training / 16\,MB validation / the
rest test; LoRA rank 8, $\alpha = 32$, batch size 1 with gradients accumulated every 8 steps;
10 epochs; learning rate $10^{-4}$; the first 10\,\% as warm-up; the base model loaded in
8-bit and trained in bfloat16.

During the reproduction we verified this conclusion in two ways.

\begin{enumerate}[label=\arabic*.,leftmargin=*]
  \item \textbf{Using our independently written program.} With the three MeDAL files, the
        same window sizes as the code path, and every parameter following the paper, differing
        only in \texttt{lora\_path}, the rate rises under all five window settings, by
        0.3987--0.4728 percentage points.
  \item \textbf{Testing the source code.} Under the 8{,}192-character window we swapped
        \texttt{lora\_path} in the authors' released program and measured the same data. The
        cross-check is as follows.
\end{enumerate}

\begin{table}[htbp]\centering\small
\caption{Cross-check with the authors' program at 8{,}192 characters.}
\begin{tabular}{@{}lrr@{}}
\toprule
Arm & Official ratio & Rate (\%) \\
\midrule
Base, no adapter & 11.204097 & 8.9253 \\
+ authors' released LoRA adapter & 10.509252 & 9.5154 \\
Difference & $-$0.694845 & +0.5901 (worse) \\
\bottomrule
\end{tabular}
\end{table}

The official code exhibits the same trend as the independent code when the LoRA adapter is
used, namely an effect in the opposite direction to the paper. The next subsection, on the
training log, shows a possible reason.

\subsubsection{Training log and checkpoint choice}

The adapter's own \texttt{trainer\_state.json} records the official training trajectory:

\begin{table}[htbp]\centering\small
\caption{Training trajectory recorded in the released adapter's \texttt{trainer\_state.json}.}
\begin{tabular}{@{}rrr@{}}
\toprule
Step & Epoch & Eval loss \\
\midrule
20 & 0.04 & 2.7491 \\
1020 & 2.0 & 2.4708 (best) \\
5120 & 10.0 & 2.6728 (degenerated by 8.2\,\%) \\
\bottomrule
\end{tabular}
\end{table}

The JSON file shows that the released adapter corresponds to the step-5120 checkpoint, whose
evaluation loss is about 8.2\,\% above the lowest value reached during training. The training
log alone is not sufficient to determine the cause of the discrepancy between the paper's
results and the results obtained from both codebases, but over-fitting during LoRA training is
a plausible contributing cause.

\subsubsection{Limitations}

The tables above show that window length affects the base model's rate, but under a given
window setting the compression penalty introduced by LoRA is larger than the difference
between adjacent window lengths.

The text tests were carried out with bfloat16 model weights and a corpus consisting of three
MeDAL files. Because the same corpus is used throughout, the results serve to compare the base
model and the adapter on the same data and under the same window conditions, but they do not
represent the overall compression rate of the full MeDAL test set.

\subsection{Audio}

The audio reproduction uses LJ001-0004.wav, 82{,}220 bytes in total.

\begin{table}[htbp]\centering\small
\caption{Audio results for the same 82{,}220-byte payload.}
\begin{tabular}{@{}lrrrr@{}}
\toprule
Implementation & Coded output & bits/byte & Rate (\%) & Ratio ($\times$) \\
\midrule
Our port (real arithmetic-coder output) & 130{,}383\,bit $=$ 16{,}298\,B & 1.5858 & 19.82 & 5.046 \\
Reference implementation (unmodified) & 16{,}399.95\,B & 1.5957 & 19.94 & 5.016 \\
Paper (LJSpeech) & --- & 1.4608 & 18.26 & 5.477 \\
\bottomrule
\end{tabular}
\end{table}

The baseline comparison is as follows.

\begin{table}[htbp]\centering\small
\caption{Audio baselines.}
\begin{tabular}{@{}lrrr@{}}
\toprule
Baseline & Rate (\%) & Paper (\%) & Difference (pp) \\
\midrule
OptimFROG \texttt{-{}-preset max} & 27.311 & 25.42 & +1.89 \\
WavPack \texttt{-hh -x3} & 35.303 & 33.15 & +2.15 \\
FLAC \texttt{-compression\_level 12} & 45.334 & 41.26 & +4.07 \\
zlib level 9 & 62.627 & --- & --- \\
\bottomrule
\end{tabular}
\end{table}

The audio reproduction succeeds. For the same 82{,}220\,B payload the independent encoder
emits 130{,}383\,bit, i.e.\ 16{,}297.875\,B, while the authors' code reports 16{,}399.95\,B;
the two rates differ by about 0.1\,\%. The authors' code result differs from the value the
paper reports on LJSpeech by about 1.68\,\%, which is likewise attributable to a difference in
sample size.

In agreement with the paper's conclusion, using a large model as the predictor for audio
markedly outperforms classical compressors, improving the rate by about 8\,\% relative to the
best classical compressor.

\subsection{Image}

The image reproduction uses five images from the CLIC professional validation set. On these
images the model generates $32\times32$ token patches sequentially along the per-channel path
supplied by the official code. In the independent implementation I use the causal distribution
at the preceding position to compute each next token's code length, in bits; per-image results
are in the table below.

\begin{table}[htbp]\centering\small
\caption{Per-image causal rates on the five CLIC images.}
\begin{tabular}{@{}lr@{}}
\toprule
Image & Causal rate (\%) \\
\midrule
alberto-montalesi-176097 & 21.4781 \\
alejandro-escamilla-6 & 17.5025 \\
ales-krivec-15949 & 18.8136 \\
alexander-shustov-73 & 41.3849 \\
allef-vinicius-109434 & 24.8664 \\
\midrule
Sample-weighted mean (5 images) & 24.9892 \\
Median & 21.4781 \\
Weighted mean excluding the highest & 20.7429 \\
\bottomrule
\end{tabular}
\end{table}

The denominator $N$ determines the amount of uncompressed data the rate is taken relative to.
In image and video processing the ambiguity in the definition of that denominator plays a
misleading role. In the authors' public code \texttt{compress\_with\_igpt.py}, the program
accumulates sample counts by height and width for each input object. The public preprocessing
function \texttt{process\_img} handles each colour channel separately: it flattens each channel
and cuts it into $32\times32$ single-channel patches, and because iGPT requires RGB input it
replicates the single-channel value into all three channels, forming a greyscale RGB input.
The ImageGPT image-processing function then quantises and tokenises this into 1{,}024 tokens.

Because each channel value is 8 bits, the accumulated sample count simultaneously represents
the number of bytes of source data actually processed, namely
$32\times32\times5733 = 5{,}870{,}592$\,B. Since one original channel sample corresponds to one
token, the token count also happens to be 5{,}870{,}592. But the tokens are the result of
iGPT's palette quantisation and are not the original channel values. This coincidence led the
original authors to believe that images and video can be reconstructed losslessly, whereas
iGPT's non-injective quantisation and the results of my decoding tests show that the process is
not lossless.

In addition, the official computation of the 15.5144\,\% rate is in error: the source code uses
the output distribution at position $t$ to encode the token at position $t$, and because that
distribution already contains information about the current token the decoder cannot obtain it
from what has already been decoded, so the requirement of causal coding is not met. That result
can be reproduced with the source code's procedure, but when the procedure is corrected to
predict the token at position $t$ from position $t-1$, the image rate becomes 23.9264\,\%.

\subsubsection{The iGPT processor and quantisation}

The ImageGPT processor reads its configuration from the checkpoint and then uses that fixed
configuration for every image; it is not retrained, the palette is not adjusted, and the
processing rules are not changed. The file \texttt{preprocessor\_config.json} supplies the
details: 512 cluster centres, with \texttt{do\_color\_quantize} and \texttt{do\_resize} both
true, so both quantisation and resizing must be enabled.

This processor is iGPT's own tokeniser, responsible for converting two-dimensional visual data
into a one-dimensional pixel sequence; iGPT is an autoregressive model over discrete colour
tokens and cannot operate without this component. The component quantises every input pixel
into a 512-colour palette, and that palette \emph{is} iGPT's vocabulary, so the model's
predictive capability cannot be used while bypassing the quantisation.

\subsubsection{Image baselines}

For the same test images and the same harness, the classical baselines are as follows.

\begin{table}[htbp]\centering\small
\caption{Image baselines on the same five CLIC images.}
\begin{tabular}{@{}lrr@{}}
\toprule
Method & Local CLIC rate (\%) & Paper CLIC2019 (\%) \\
\midrule
WebP (lossless) & 39.362 & 35.56 \\
JPEG-2000 (lossless) & 50.215 & 48.64 \\
PNG (ffmpeg, compression level 100) & 76.848 & 47.08 \\
\bottomrule
\end{tabular}
\end{table}

Overall the baselines are reproduced in reasonable agreement with the paper's data, and the
compression rate achieved by the large model is markedly higher than the classical baselines.
But because the images processed by iGPT do not undergo a lossless procedure, this comparison
is fundamentally unfair.

\subsubsection{The premise of per-channel processing}

As noted earlier, the original paper does not feed a mixed RGB three-channel input into iGPT;
it encodes each channel independently. The paper's stated reason is:

\begin{quote}
``Because the tokens of iGPT are greyscale values, we regard an image as three greyscale
images, one for each RGB channel \ldots{} During decompression, each channel is decompressed
independently, and the three channels are then merged to reconstruct the original image.''
\end{quote}

I checked this premise against the checkpoint's own palette and found it to be incorrect. The
512 cluster centres of iGPT are in fact the result of quantising full RGB colours: all 512
palette entries are RGB triples, and \emph{none} of them is a grey colour. Even under the
near-grey criterion set out in the paper, only 61 are close to grey. This shows that iGPT's
tokens are not greyscale values; the per-channel method does not introduce serious error, but
adopting it is not necessary. The check is tabulated below.

\begin{table}[htbp]\centering\small
\caption{Inspection of the iGPT palette.}
\begin{tabular}{@{}lr@{}}
\toprule
Check & Result \\
\midrule
Shape of \texttt{clusters} & $(512, 3)$ --- 512 RGB triples \\
Entries exactly grey ($R=G=B$) & 0 / 512 \\
Near grey (maximum channel deviation $<0.05$) & 61 / 512 \\
\bottomrule
\end{tabular}
\end{table}

\subsubsection{Reversibility of the palette mapping}

Can the original channel value be recovered uniquely from the token alone, once it has passed
through the iGPT processor? This subsection demonstrates empirically that it cannot, by
enumerating all possible inputs and by inspecting real image patches.

The image loader first splits the image by colour channel, so \texttt{language\_model.compress}
receives a tensor of shape $(32,32,1)$. Because the last dimension is 1, the code replicates
that channel into the three RGB positions and hands a tensor of shape $(32,32,3)$ to the
processor. Every 8-bit channel value $v$ therefore becomes the greyscale triple $(v,v,v)$,
which is then mapped to a token ID; the encoding stores only that token ID and does not retain
the original channel value.

For a greyscale 8-bit channel the input range is $(0,255)$, i.e.\ 256 values, yet after palette
quantisation these 256 values use only \textbf{48} of the 512 tokens. Next, in a value-by-value
enumeration the program groups by token ID and counts collisions: if a token receives $k$ input
values then those inputs form $\binom{k}{2}$ colliding pairs. Over the 48 tokens there are
\textbf{585} colliding pairs in total. The largest groups---tokens 391, 324, 500, 389 and
48---each contain 7 input values. The same many-to-one mapping occurs on real images.

The examples above show that iGPT's palette quantisation is not injective, so the input value
cannot be recovered uniquely from the output token, and the present compression pipeline
transmits no additional residual information that would let the decoder decode losslessly.

\subsection{Video}

\subsubsection{iGPT results}

\begin{table}[htbp]\centering\small
\caption{Video results.}
\begin{tabular}{@{}lr@{}}
\toprule
Quantity & Value \\
\midrule
Coverage of the iGPT measurement & 9 distinct frames \\
Processing units per frame & 25{,}920 single-channel patches \\
Samples per frame counted in the denominator & 26{,}542{,}080 \\
Per-channel causal rate & 26.1866\,\% \\
Paper SI Table 4 high-resolution result & 15.31\,\% (locally 10.88\,pp higher) \\
H.265 baseline on the same clip & 26.467\,\% (numerically 0.2804\,pp higher; the objects coded differ, so no same-basis verdict) \\
Object coded & $32\times32$, 512-colour palette token representation; not all source samples recoverable \\
\bottomrule
\end{tabular}
\end{table}

The iGPT per-channel measurement covers 9 distinct frames and gives a sample-weighted rate of
26.1866\,\%, with a frame-to-frame range of 0.0316 percentage points. The paper's result is
15.31\,\%, about 10.88 percentage points lower than the independent measurement. The cause of
this gap is the same as for the image modality: the official public code shares
\texttt{compress\_with\_igpt.py} and \texttt{language\_model.compress()} between images and
video, so the conclusion that using position-$t$ information to estimate the code length of the
token at position $t$ makes the computed result optimistic applies to video as well.

Moreover, since both image and video use iGPT as the predictor, the conclusion that palette
quantisation makes the compression procedure non-lossless also applies to the video modality.

\subsubsection{Classical video compression baselines}

The classical video baselines cover 10 frames of the Netflix Aerial 4K clip, whose source format
is $4096\times2160$, 60\,fps, YUV420. The results are as follows.

\begin{table}[htbp]\centering\small
\caption{Classical video baselines on the same clip.}
\begin{tabular}{@{}lrrrr@{}}
\toprule
Encoder & Compressed size (B) & Rate (\%) & Time (s) & Paper (\%) \\
\midrule
H.265 crf 0 (x265 veryslow) & 35{,}124{,}565 & 26.467 & 154.4 & 21.69 \\
H.264 crf 0 (x264 veryslow) & 45{,}236{,}644 & 34.087 & 5.5 & 37.42 \\
FFV1 level 3 & 53{,}234{,}817 & 40.114 & 1.8 & 40.16 \\
zstd \texttt{-{}-ultra -22} (raw YUV) & 61{,}434{,}434 & 46.292 & 119.1 & 51.04 \\
\bottomrule
\end{tabular}
\end{table}

The classical baselines are reproduced essentially successfully, with rate differences
fluctuating between 0.06\,\% and 4.7\,\%.

\appendix
\section{Translator's notes: discrepancies found in the source}

The following are internal inconsistencies in \texttt{CN.docx}. They are recorded here rather
than silently corrected, because the source was not to be modified.

\begin{enumerate}[label=\arabic*.,leftmargin=*]
  \item \textbf{Table 5 (``Text rates by window unit and length'') is inconsistent with
        Table 6 and with the running text.} Table 5 gives the base rates as 6.5322 / 6.3927 /
        6.3321 / 6.1657 / 6.1023 and the penalties as +0.3997 to +0.4728, whereas Table 6 gives
        the base range as 8.8038--9.4240\,\% and the penalties as +0.5752 to +0.6821. The factor
        between the two sets is $\log_2 e = 1.4427$: the Table 5 values are natural-log code
        lengths (``nats'') rather than bits. One of the two tables should be brought into line
        with the other.
  \item \textbf{Section 4.2, item 1} states the penalty as ``0.3987--0.4728 percentage points'',
        which repeats the un-converted values, and so conflicts with Table 6 in the same
        section.
  \item \textbf{Section 4.2} says the 8{,}192-character result ``differs by about 0.05
        percentage points'' from the paper's 9.08\,\%, whereas the abstract and Section 4.1 give
        +0.08\,pp against 9.06\,\%. Both comparisons are defensible (9.08 is the Si Table 6
        base row, 9.06 the Si Table 9 context-8{,}192 row) but the report should say which row
        is being compared.
  \item \textbf{Section 3.6} says bGPT uses ``19 bytes per patch'' while returning ``17 slot
        predictions''. A 17-slot row implies a 16-byte payload with one additional predicted
        byte, so 19 appears to be a typo for 17, or the patch definition needs restating.
  \item \textbf{Section 4.4.4} writes the corrected image rate as ``23. 9264\,\%'' with a space
        inside the number.
  \item \textbf{Section 4.4.3} says ``about seven percentage points'' from the paper's
        14.62\,\% for a median of 21.4781\,\%; the difference is 6.86\,pp, and the weighted
        mean is 10.37\,pp. Stating which statistic is meant would remove the ambiguity.
\end{enumerate}

\end{document}
